\chapter{Plateformes centralisées}\label{ch:m3:centralised-exchanges}

Durant la semaine commencée le 6 novembre 2022, les clients de l'une des plus grandes plateformes de cryptomonnaies ont retiré
sept milliards de dollars de celle-ci. Le 8 novembre, elle a suspendu tous les retraits de ses clients~; le 11 novembre, elle a déposé
son bilan. Ses clients la considéraient comme une plateforme de négociation~: un lieu où
leurs ordres rencontraient d'autres ordres. Ils ont appris qu'elle avait aussi été leur courtier, tenant leurs comptes~;
leur dépositaire, conservant leurs coins~; leur chambre de compensation, garantissant leurs transactions et liquidant
leurs positions à effet de levier~; et, par l'intermédiaire d'une société de trading affiliée qui pouvait emprunter auprès de la plateforme
sans limite, la contrepartie qui avait dépensé leur argent. La négociation de cryptomonnaies se déroule surtout sur de telles
plateformes. Ce chapitre décrit ce qu'est une \omterm{def:m3:centralised-exchanges:cex}{plateforme centralisée}, comment une société de trading s'y connecte et en paie
le prix, et à quoi elle est exposée tant que ses actifs y demeurent~: le risque que la plateforme ne soit pas ce qu'elle prétend être
ce que les \omterm{def:m3:centralised-exchanges:por}{preuves de réserves} peuvent et ne peuvent pas établir, et les dispositifs de règlement
conçus pour que le collatéral d'une société reste hors de portée de la plateforme.

\section{La plateforme qui est aussi courtier, dépositaire et chambre de compensation}

Sur les marchés d'actions et de contrats à terme du livre~1, une transaction traverse des institutions distinctes,
chacune réglementée pour son rôle~: un courtier tient le compte du client et achemine ses ordres (chapitre~4), la bourse les apparie,
une contrepartie centrale garantit la transaction et perçoit les appels de marge, et un dépositaire ou un système de règlement-livraison
détient les titres et les espèces (chapitre~5). Chacun possède son propre capital, son propre superviseur et ses propres règles
sur ce qu'il peut faire des actifs de ses clients.

\begin{definition}[Plateforme centralisée]\label{def:m3:centralised-exchanges:cex}
Une \emph{plateforme centralisée}\index{plateforme centralisée} (centralised exchange) est une plateforme de négociation de cryptomonnaies exploitée par une société
qui détient les actifs de ses clients dans des portefeuilles et des comptes bancaires qu'elle contrôle, tient leurs soldes dans son propre
registre interne, apparie leurs ordres dans son propre carnet d'ordres et règle les transactions en modifiant ce registre
plutôt que sur une \omterm{def:m3:blockchains-for-traders:chain}{blockchain}.
\end{definition}

Seuls les dépôts et les retraits touchent une chaîne. Un client qui envoie des bitcoins à la plateforme abandonne ses coins en échange
d'une ligne dans la base de données de la plateforme~; chaque transaction ensuite déplace des montants d'une ligne à l'autre, et seul
un retrait transforme de nouveau une ligne en coins. La même société assure les quatre rôles de la chaîne traditionnelle, et
généralement davantage~: elle prête à ses clients qui négocient avec effet de levier, exploite le moteur de liquidation qui
clôt leurs positions (\cref{ch:m3:margin-liquidation-and-loss-allocation}), émet son propre \omterm{def:m3:blockchains-for-traders:key}{jeton} et peut exploiter un desk de trading
ou être affiliée à un tel desk.

\begin{omfigure}
\centering
\begin{tikzpicture}[font=\small,
  r/.style={draw,thick,rounded corners=2pt,align=center,minimum height=0.8cm,text width=2.3cm},
  arr/.style={-{Stealth[length=2.2mm]},thick}]
\node[font=\small\bfseries,anchor=west] at (-1.3,2.1) {Titres et contrats à terme};
\node[r,draw=omDef,fill=omDef!8] (b) at (0,1.2) {courtier\\(comptes)};
\node[r,draw=omThm,fill=omThm!8] (x) at (3.1,1.2) {bourse\\(appariement)};
\node[r,draw=omProp,fill=omProp!8] (c) at (6.2,1.2) {contrepartie\\centrale};
\node[r,draw=omExo,fill=omExo!8] (k) at (9.3,1.2) {dépositaire,\\règlement-livraison};
\draw[arr] (b) -- (x); \draw[arr] (x) -- (c); \draw[arr] (c) -- (k);
\node[font=\small\bfseries,anchor=west] at (-1.3,-0.3) {Plateforme centralisée de crypto-actifs};
\node[draw=omThm,thick,rounded corners=3pt,fill=omThm!4,minimum width=10.9cm,minimum height=1.2cm,
  align=center] (one) at (4.65,-1.2) {une seule société~: comptes, appariement, marge et liquidation,\\
  garde des coins et des espèces, prêt, jeton propre, (souvent) un desk de trading affilié};
\end{tikzpicture}

\omcaption{Sur les marchés réglementés de titres et de contrats à terme, des institutions distinctes tiennent le compte, apparient
l'ordre, garantissent la transaction et détiennent les actifs~; une \omterm{def:m3:centralised-exchanges:cex}{plateforme centralisée} de crypto-actifs fait tout cela au sein
d'une seule société, dont le registre interne est la seule trace de qui possède quoi. Schéma.}\label{fig:m3:centralised-exchanges:roles}
\end{omfigure}

\begin{definition}[Ségrégation des actifs des clients]\label{def:m3:centralised-exchanges:segregation}
La \emph{ségrégation des actifs des clients}\index{ségrégation des actifs des clients} (customer asset segregation) est la séparation des actifs détenus
pour le compte des clients d'avec les actifs propres du détenteur, dans des comptes ou des portefeuilles distincts et selon des modalités juridiques qui
les tiennent à l'écart du patrimoine du détenteur en cas de défaillance, de sorte qu'ils ne puissent pas servir à ses propres fins et
reviennent aux clients en cas d'insolvabilité.
\end{definition}

La ségrégation est ce qui rend supportable pour ses clients la défaillance d'un courtier ou d'un futures
commission merchant. Sur une \omterm{def:m3:centralised-exchanges:cex}{plateforme centralisée}, elle dépend des conditions générales de la plateforme, du droit du pays où elle est
immatriculée, et de son honnêteté. Lorsque les conditions ou la loi sont muettes, le solde d'un client constitue une créance
non garantie sur la société, et les coins que détient la plateforme appartiennent à l'ensemble
de ses créanciers.

Pour une société de trading, la conséquence pratique est qu'un solde sur une plateforme constitue une exposition de crédit
envers cette plateforme, pour le montant total, aussi longtemps qu'il y demeure. La société fixe des limites par plateforme, comme une
banque fixe des limites par contrepartie, et mesure ce qu'elle perdrait, ainsi que la rapidité avec
laquelle elle pourrait retirer ses fonds, si la plateforme cessait de payer.

\section{Comptes, appariement et interfaces}

Une société ouvre un compte maître auprès de la plateforme après un examen d'entrée en relation (\cref{ch:m3:getting-access-crypto-venues})
et le divise selon ses stratégies et ses desks.

\begin{definition}[Sous-compte, clé d'API]\label{def:m3:centralised-exchanges:subaccount}
Un \emph{sous-compte}\index{sous-compte} (sub-account) est un compte margé séparément sous un compte maître, avec
ses propres soldes, positions et limites, entre lesquels le compte maître peut transférer des actifs. Une
\emph{clé d'API}\index{clé d'API} (API key) est un identifiant, composé d'un identifiant de clé et d'un secret, avec lequel un programme signe ses
requêtes adressées à une plateforme~; chaque clé porte des permissions (lecture des soldes, négociation, retrait) et peut être
restreinte à une liste d'adresses IP.
\end{definition}

Les \omterm{def:m3:centralised-exchanges:subaccount}{sous-comptes} séparent les stratégies, de sorte qu'une perte ou une liquidation dans l'une ne puisse pas consommer le
collatéral d'une autre, et ils permettent de suivre les frais, les positions et le P\&L de chaque stratégie sans comptabilité propre à la
société. Les \omterm{def:m3:centralised-exchanges:subaccount}{clés d'API} sont, après les \omterm{def:m3:blockchains-for-traders:key}{clés privées} de ses portefeuilles, le secret opérationnel le plus sensible de la
société~: une clé dotée d'un droit de retrait est une clé de l'argent. Les clés de négociation sont émises sans droit de retrait,
restreintes aux adresses des serveurs de la société, renouvelées et conservées dans un coffre-fort de secrets plutôt que dans le code~;
les retraits passent par un processus distinct, avec ses propres approbations et ses listes blanches d'adresses.

L'appariement suit la priorité prix-temps du livre~1, chapitre~19, avec des différences qu'un trader
remarque aussitôt. Le marché ne ferme jamais, de sorte qu'il n'y a ni enchères d'ouverture ni enchères de clôture, ni écart d'une nuit
à l'autre~; les types d'ordres sont moins nombreux, avec les ordres au marché, à cours limité et stop, leurs indicateurs de durée de validité, et
l'\omterm{def:m3:centralised-exchanges:postonly}{ordre post-only} de la section suivante~; et la prévention de l'auto-négociation, qui empêche les ordres d'achat et de vente d'une même société de
s'apparier entre eux, est un indicateur attaché à chaque ordre, avec plusieurs modes (annuler l'ordre au repos,
l'ordre entrant, ou les deux).

Un programme communique avec une plateforme par deux types d'interfaces. Les requêtes (passer, modifier et annuler un ordre,
interroger les soldes et les ordres ouverts) transitent par une interface requête-réponse, généralement en HTTPS~; les données de
marché et les événements propres au compte (exécutions, variations de solde) arrivent par des connexions WebSocket en flux continu. Le
comportement de ces flux, leurs numéros de séquence, leurs instantanés et leurs lacunes fait l'objet du
\cref{ch:m3:crypto-data-and-infrastructure}. Le côté requêtes est celui où la plateforme se protège.

\begin{definition}[Limite de débit, poids de requête]\label{def:m3:centralised-exchanges:ratelimit}
Une \emph{limite de débit}\index{limite de débit} (rate limit) est le nombre maximal de requêtes, ou d'unités de charge, que
une plateforme accepte d'un client dans un intervalle de temps donné, au-delà duquel elle rejette les requêtes et peut bannir le client
pendant une période. Un \emph{poids de requête}\index{poids de requête} (request weight) est le nombre d'unités qu'une requête consomme sur une telle limite,
fixé par la plateforme pour chaque point d'accès afin de refléter son coût~: un instantané complet
du carnet d'ordres pèse davantage qu'un ordre unique.
\end{definition}

\begin{dated}{2026-09}{Limites de débit sur une grande plateforme}\label{dat:m3:centralised-exchanges:limits}
La documentation de l'API spot de Binance attribue à chaque point d'accès un poids et renvoie, dans un en-tête de réponse, le poids consommé jusqu'ici
dans chaque intervalle~; les limites s'appliquent par adresse IP, et non par \omterm{def:m3:centralised-exchanges:subaccount}{clé d'API}, et les limites de nombre d'ordres
par compte, également signalées dans un en-tête. Une requête qui dépasse une limite reçoit le statut HTTP 429
avec un en-tête \texttt{Retry-After}~; une adresse IP qui continue d'émettre après des réponses 429 est automatiquement
bannie, avec le statut 418, pour des durées qui passent, pour les récidivistes, de 2 minutes à 3 jours. Le 24
septembre 2026, son point d'accès d'information sur l'échange renvoyait des limites de 6\,000 unités de \omterm{def:m3:centralised-exchanges:ratelimit}{poids de requête} par
minute, 100 ordres toutes les 10 secondes, 200\,000 ordres par jour et 300\,000 requêtes brutes toutes les 5 minutes~;
ces valeurs changent, et un client les lit plutôt que de les coder en dur.
\end{dated}

Deux propriétés de ces limites conditionnent la conception d'un client. La première est que les fenêtres sont fixes.

\begin{proposition}[La rafale à la frontière d'une fenêtre]\label{prop:m3:centralised-exchanges:burst}
Avec une limite de $L$ unités par fenêtre fixe de longueur $T$ alignée sur l'horloge, un client peut émettre $2L$
unités dans un intervalle de longueur inférieure à $T$ à cheval sur une frontière de fenêtre, mais jamais plus de
$2L$ dans un intervalle quelconque de longueur $T$. Une fenêtre glissante, qui compte les unités des $T$
dernières secondes, borne chaque intervalle de longueur $T$ à $L$.
\end{proposition}

\begin{proof}
Émettre $L$ unités juste avant la frontière et $L$ juste après~: chaque fenêtre contient $L$. Tout intervalle de
longueur $T$ rencontre au plus deux fenêtres, dont chacune contient au plus $L$. Une fenêtre glissante compte, à
chaque instant, exactement les unités émises dans l'intervalle de longueur $T$ qui précède et refuse toute unité qui
la ferait dépasser $L$.
\end{proof}

Un client qui s'appuie sur la rafale à la frontière en tire un avantage l'espace d'un instant, et se fait parfois
prendre par une plateforme qui compte autrement que ce qu'elle documente. La conception robuste compte de façon
prudente de son côté, n'envoie jamais une requête dont elle prévoit le refus, et traite un 429 comme un bogue à
corriger, car le bannissement qui suit des refus répétés éloigne la société de trading de la plateforme.

La seconde propriété compte davantage pour un teneur de marché. Toute modification de cotation est un ordre, ou une
annulation suivie d'un ordre, et le nombre d'ordres est limité par compte. Lorsque le marché bouge vite, la stratégie
veut modifier ses cotations le plus souvent, précisément quand la limite est atteinte. Un client qui met en file les
mises à jour qu'il ne peut pas encore envoyer les envoie plus tard, dans l'ordre~: des cotations calculées pour un
marché qui a bougé entre-temps. \Cref{fig:m3:centralised-exchanges:burst} montre l'effet.

\begin{omfigure}
\begin{tikzpicture}
\begin{axis}[width=11cm,height=5cm,axis y line*=left,
  xlabel={seconde},ylabel={mises à jour de cotation en file},
  tick label style={font=\small},label style={font=\small},
  xmin=0,xmax=299,ymin=0,ymax=1000,grid=major,grid style={gray!25},
  table/col sep=comma]
\addplot[omThm,very thick] table[x=second,y=fifo]{figdata/markets-3/15-centralised-exchanges/burst.csv};\label{plot:m3:burstfifo}
\addplot[omExo,very thick,dashed] table[x=second,y=coalesced]{figdata/markets-3/15-centralised-exchanges/burst.csv};\label{plot:m3:burstcoal}
\end{axis}
\begin{axis}[width=11cm,height=5cm,axis y line*=right,axis x line=none,
  ylabel={mise à jour la plus ancienne, s},ylabel near ticks,
  tick label style={font=\small},label style={font=\small},
  xmin=0,xmax=299,ymin=0,ymax=100,
  legend columns=3,legend style={font=\footnotesize,at={(0.5,-0.3)},anchor=north,draw=none,fill=none,/tikz/every even column/.append style={column sep=6pt}},
  table/col sep=comma]
\addlegendimage{/pgfplots/refstyle=plot:m3:burstfifo}\addlegendentry{en file, envoyées dans l'ordre (gauche)}
\addlegendimage{/pgfplots/refstyle=plot:m3:burstcoal}\addlegendentry{regroupées (gauche)}
\addplot[omDef,thick] table[x=second,y=age]{figdata/markets-3/15-centralised-exchanges/burst.csv};\addlegendentry{âge de la plus ancienne (droite)}
\end{axis}
\end{tikzpicture}

\omcaption{Une stratégie de cotation sur cinq instruments sous une limite de 100 ordres toutes les 10 secondes~: 5 mises
à jour de cotation par seconde, 25 par seconde pendant une minute agitée (secondes 60 à 119). En file et envoyées dans l'ordre, l'arriéré
atteint 900 mises à jour et la plus ancienne attend 89 secondes~; regroupées (seule la dernière cotation conservée par
instrument), rien n'attend. Taux illustratifs~; les limites sont celles de l'\cref{dat:m3:centralised-exchanges:limits}.
Données~: le tutoriel du chapitre.}\label{fig:m3:centralised-exchanges:burst}
\end{omfigure}

Le remède consiste à regrouper~: conserver, par instrument et par côté, seule la dernière cotation souhaitée, et envoyer
l'écart entre celle-ci et ce qui repose dans le carnet lorsque la limite le permet. L'arriéré est alors d'au plus une
mise à jour par instrument et par côté, et ce qui est envoyé est à jour. La construction de ce chapitre est un limiteur
de débit qui décide, avant chaque requête, si la plateforme l'acceptera.

\section{Frais, rabais et programmes de teneurs de marché}

Les plateformes de crypto-actifs facturent selon une tarification maker-taker avec des paliers de volume (livre~1, chapitre~4)~: un preneur,
dont l'ordre s'exécute à l'arrivée, paie davantage qu'un donneur, dont l'ordre au repos fournit la liquidité, et les deux
paient moins à mesure que leur volume sur trente jours augmente. Une stratégie qui entend fournir de la liquidité doit s'assurer
que ses ordres ne la prennent jamais par accident.

\begin{definition}[Ordre post-only]\label{def:m3:centralised-exchanges:postonly}
Un \emph{ordre post-only}\index{ordre post-only} est un ordre à cours limité que la plateforme n'accepte que s'il
reposerait dans le carnet~; s'il devait s'exécuter aussitôt contre un ordre au repos, la plateforme le rejette (ou, sur
certaines plateformes, le reprice d'un tick) au lieu de l'exécuter comme preneur.
\end{definition}

Un \omterm{def:m3:centralised-exchanges:postonly}{ordre post-only} garantit les frais de donneur et l'absence de transaction croisée, au prix d'un rejet lorsque le
marché a franchi le cours de l'ordre pendant son trajet. Pour un teneur de marché dont les cotations sont calculées à
partir d'un prix vieux de quelques millisecondes, ce rejet est le bon résultat.

\begin{dated}{2026-09}{Grilles de frais spot de deux plateformes}\label{dat:m3:centralised-exchanges:fees}
La grille de frais spot de Binance facture à un utilisateur ordinaire 0.100~\% en donneur comme en preneur (0.075~\% lorsqu'il paie
avec son propre \omterm{def:m3:blockchains-for-traders:key}{jeton}, le BNB), et son palier le plus élevé, VIP~9 (volume sur trente jours d'au moins 4~milliards d'USD
ou une détention de BNB), 0.011~\% en donneur et 0.023~\% en preneur. Kraken Pro facture son palier spot le plus bas 0.40~\% en donneur
et 0.80~\% en preneur, et son palier le plus élevé 0.00~\% et 0.05~\%. Les grilles changent souvent, et les grands teneurs
de marché négocient des conditions en dehors d'elles.
\end{dated}

\begin{table}[ht]
\centering
\small
\begin{tabular}{@{}lrrrr@{}}
\toprule
Plateforme et palier & maker (pb) & taker (pb) & 2 jambes taker (pb) & 2 jambes maker (pb) \\
\midrule
Binance, ordinaire & 10.0 & 10.0 & 20.0 & 20.0 \\
Binance, VIP 9 & 1.1 & 2.3 & 4.6 & 2.2 \\
Kraken Pro, plus bas & 40.0 & 80.0 & 160.0 & 80.0 \\
Kraken Pro, plus élevé & 0.0 & 5.0 & 10.0 & 0.0 \\
\bottomrule
\end{tabular}
\caption{Le coût d'un aller-retour (un achat et une vente) en points de base de notionnel, d'après les grilles de
l'\cref{dat:m3:centralised-exchanges:fees}. Une stratégie dont l'edge attendu par aller-retour se chiffre en quelques
points de base n'existe que dans les paliers supérieurs.}\label{tab:m3:centralised-exchanges:fees}
\end{table}

\Cref{tab:m3:centralised-exchanges:fees} illustre bien l'effet des paliers. Une stratégie qui capte cinq
points de base par aller-retour perd de l'argent au palier le plus bas de l'une ou l'autre plateforme et en gagne au palier le plus élevé~;
la différence ne tient pas à l'habileté mais au volume. Les programmes de teneur de marché des plateformes vont plus loin, en versant des rétrocessions
(frais maker négatifs) et parfois un montant fixe, en contrepartie d'obligations de cotation~: une fourchette maximale,
une taille minimale, une part du temps. Ce sont des programmes d'incitation au sens du livre~1, et
\cref{ch:m3:getting-access-crypto-venues} en détaille un exemple. Deux mises en garde s'imposent. Une remise versée dans le
\omterm{def:m3:blockchains-for-traders:key}{jeton} propre à la plateforme est une position sur ce \omterm{def:m3:blockchains-for-traders:key}{jeton}, dont la valeur dépend de la santé de la plateforme. Et une
grille fondée sur le volume récompense le volume, y compris celui qui s'échange contre lui-même, ce qui explique en partie pourquoi
les volumes publiés par certaines plateformes surestiment l'activité réelle (du \cref{ch:m3:spot-markets}).

\section{La plateforme comme contrepartie~: la faillite de FTX}

Les clients de FTX ont découvert en une semaine la différence entre une plateforme et une contrepartie. Le récit
qui suit provient des poursuites engagées par la SEC américaine et des pièces
transmises par les débiteurs, c'est-à-dire les administrateurs judiciaires qui ont repris les sociétés en faillite.

FTX exploitait une grande plateforme internationale, FTX.com, et une plus petite plateforme américaine, FTX.US. Alameda Research, une
société de trading appartenant au même fondateur, est parvenue à dominer les emprunts sur FTX.com. En décembre 2022
la SEC a accusé le fondateur, Samuel Bankman-Fried, de fraude au détriment des investisseurs au capital auprès desquels FTX avait
levé plus de 1.8~milliard de dollars, en dissimulant que les fonds des clients de FTX étaient détournés vers Alameda,
et que les fonds clients, confondus avec les siens, servaient à des investissements de capital-risque, à l'immobilier et à des dons politiques.
Les rapports des débiteurs en décrivent le mécanisme. Le code de la plateforme accordait à Alameda une limite d'emprunt de 65~
milliards de dollars et des paramètres lui permettant de retirer des actifs et de détenir un solde négatif sans être liquidé,
privilèges dont aucun autre client ne bénéficiait. Les dépôts des clients en dollars ont été en partie reçus sur des comptes
contrôlés par Alameda. Presque tous les actifs cryptographiques étaient conservés dans des portefeuilles connectés à internet, sans
les contrôles multisignatures courants dans le secteur.

\begin{omfigure}
\begin{tikzpicture}
\begin{axis}[width=11cm,height=5cm,ybar,bar width=9pt,
  xlabel={novembre 2022},ylabel={millions de dollars},
  tick label style={font=\small},label style={font=\small},
  xmin=0.4,xmax=11.6,xtick={1,2,3,4,5,6,7,8,9,10,11},ymin=-500,ymax=8000,grid=major,grid style={gray!25},
  legend columns=3,legend style={font=\footnotesize,at={(0.5,-0.3)},anchor=north,draw=none,fill=none,/tikz/every even column/.append style={column sep=10pt}},
  table/col sep=comma]
\addplot[fill=omThm!60,draw=omThm,area legend] table[x=day,y=cust_out]{figdata/markets-3/15-centralised-exchanges/flows.csv};\addlegendentry{retraits nets des clients}
\addplot[omDef,very thick,sharp plot,mark=*,mark size=1.5pt] table[x=day,y=cust_cum]{figdata/markets-3/15-centralised-exchanges/flows.csv};\addlegendentry{cumulé}
\addplot[omExo,very thick,sharp plot,dashed,mark=square*,mark size=1.5pt] table[x=day,y=rel_cum]{figdata/markets-3/15-centralised-exchanges/flows.csv};\addlegendentry{dépôts nets des parties liées, cumulés}
\end{axis}
\end{tikzpicture}

\omcaption{FTX.com, du 1\up{er} au 11 novembre 2022~: les clients ont retiré, net, 7.0~milliards de dollars, dont 3.3~milliards
le 7 novembre~; Alameda et d'autres parties liées ont déposé, net, 3.2~milliards de dollars, surtout du 6 au 8 novembre.
En millions de dollars, aux cours de la date de la requête, tels que publiés~; les sommes cumulées diffèrent des totaux du dossier par
des arrondis. Données~: FTX Debtors, présentation du 2 mars 2023, In re FTX Trading Ltd., Bankr.\ D.\ Del.\ No.\
22-11068.}\label{fig:m3:centralised-exchanges:flows}
\end{omfigure}

Quand les doutes sur le bilan d'Alameda ont été rendus publics au début de novembre 2022, les clients se sont précipités vers les sorties.
\Cref{fig:m3:centralised-exchanges:flows} montre cette ruée dans les chiffres des débiteurs~: retraits nets de clients sur FTX.com,
d'environ 1.8~milliard de dollars le 6 novembre et 3.3~milliards le 7 novembre, avant que les retraits ne cessent de fait.
Les sociétés ont déposé leur bilan au titre du Chapter~11 le 11 novembre (et certaines filiales le 14 novembre).
Le jour du dépôt, des transferts non autorisés ont encore vidé les portefeuilles des plateformes de quelques centaines de
millions de dollars.

\begin{table}[ht]
\centering
\small
\begin{tabular}{@{}lrrr@{}}
\toprule
FTX.com, millions de dollars & créances des clients & actifs localisés & couverture \\
\midrule
Espèces et \omterm{def:m3:blockchains-for-traders:stable}{stablecoins} & 6\,991 & 270 & 3.9~\% \\
Bitcoin & 1\,591 & 1 & 0.1~\% \\
Ether & 922 & 9 & 1.0~\% \\
Ensemble des \omterm{def:m3:blockchains-for-traders:key}{jetons} liquides (catégorie A) & 10\,544 & 694 & 6.6~\% \\
Ensemble des \omterm{def:m3:blockchains-for-traders:key}{jetons}, créances comprises & 11\,233 & 2\,540 & 22.6~\% \\
\bottomrule
\end{tabular}
\caption{Créances des clients sur FTX.com et actifs localisés par les débiteurs dans ses portefeuilles à la date de la requête,
en millions de dollars, aux cours de la date de la requête. Catégorie A~: \omterm{def:m3:blockchains-for-traders:key}{jetons} dont la capitalisation boursière atteint au moins 15~
millions de dollars et dont le volume quotidien atteint au moins 1~million de dollars~; la plupart des actifs de catégorie B étaient des \omterm{def:m3:blockchains-for-traders:key}{jetons} liés à FTX et à
Alameda. Les créances sont ce que les clients devaient à la plateforme. Chiffres préliminaires, tirés de la présentation des débiteurs
du 2 mars 2023.}\label{tab:m3:centralised-exchanges:shortfall}
\end{table}

\Cref{tab:m3:centralised-exchanges:shortfall} est le bilan auquel les clients étaient en réalité exposés. Sur
10.5~milliards de dollars de créances sur des actifs liquides, la plateforme en détenait 694 millions. Elle détenait un million de dollars
de bitcoins pour 1.6~milliard de dollars dus. Les seuls actifs dont la valeur dépassait celle des créances qui s'y rattachaient étaient
des \omterm{def:m3:blockchains-for-traders:key}{jetons} peu liquides liés au groupe lui-même, valorisés à des cours auxquels personne ne pouvait vendre. La plainte ultérieure
des débiteurs contre les fondateurs relève qu'en août 2022 ils estimaient en privé que la plateforme devait à ses clients plus de
8~milliards de dollars qu'elle ne pouvait pas rembourser.

\begin{dated}{2026-09}{FTX après coup}\label{dat:m3:centralised-exchanges:ftx}
Samuel Bankman-Fried a été reconnu coupable de fraude par un jury de New York en novembre 2023 et condamné le 28
mars 2024 à 25 ans de prison, avec une confiscation de plus de 11~milliards de dollars. Le 8 août 2024, un tribunal fédéral,
saisi dans une affaire engagée par la CFTC, a ordonné à FTX et à Alameda de
verser 8.7~milliards de dollars de restitutions et 4~milliards au titre des gains illicites, au motif que FTX avait
prétendu séparer les actifs des clients tout en les confondant et en les détournant. Selon le US Bankruptcy Code, une créance
est arrêtée en dollars à la date de la requête. Un client à qui un bitcoin était dû était donc créancier à hauteur de sa valeur en dollars de novembre
2022~: le plan confirmé le 7 octobre 2024 promettait à 98~\% des créanciers, en nombre, environ 119~\% de
leurs créances admises, en dollars, quelle qu'ait été ensuite la valeur des pièces.
\end{dated}

FTX n'a pas été la première faillite de ce type. Mt.~Gox, une des premières plateformes de bitcoin, a déposé son bilan à Tokyo
le 28 février 2014, en indiquant qu'elle avait pu perdre environ 750\,000 bitcoins de ses clients et 100\,000 de
ses propres bitcoins~; la procédure a été convertie en réhabilitation civile en juin 2018, et
l'administrateur judiciaire a commencé à rembourser les créanciers en bitcoins et en bitcoin cash en 2024, dix ans après la faillite. Les
leçons pour une société de trading sont très pratiques. Un solde détenu sur une plateforme est un prêt non garanti consenti à celle-ci, à moins que la société
n'en sache autrement. Ce prêt doit être dimensionné comme tel, réparti entre plusieurs plateformes et rapatrié
vers la conservation propre de la société dès qu'il n'est plus nécessaire au trading. Une ruée commence avant que l'information ne soit certaine, et les premiers
à retirer sont payés. Et les déclarations d'une plateforme sur ses réserves constituent une créance, que la section suivante
examine.

\section{Preuve de réserves et ses limites}

Après novembre 2022, plusieurs plateformes ont commencé à publier des éléments attestant qu'elles détenaient les actifs de leurs clients.

\begin{definition}[Preuve de réserves, preuve de passif]\label{def:m3:centralised-exchanges:por}
Une \emph{preuve de réserves}\index{preuve de réserves} est un élément montrant qu'une plateforme contrôle des actifs on-chain
au moins égaux aux soldes de ses clients à un instant donné~: ce sont des signatures de ses adresses, dont les soldes
que chacun peut lire sur les chaînes, comparées au total des soldes clients. Une \emph{preuve de
passif}\index{preuve de passif} est la preuve que ce total est complet et exact~: un engagement sur chaque solde client,
généralement la racine d'un arbre de Merkle (un arbre binaire dans lequel chaque nœud
est le hachage de ses fils), grâce à quoi chaque client peut vérifier que son propre solde a bien été compté.
\end{definition}

Dans un arbre de Merkle à sommes, chaque feuille engage un compte, en hachant l'identifiant du compte, un nonce secret
connu de son seul titulaire, et son solde, et elle porte ce solde~; chaque nœud parent hache ensemble ses deux fils
avec la somme de leurs soldes, et porte cette somme. Le hachage de la racine engage chaque feuille, et sa somme est
le total publié des passifs. Un client reçoit le nonce de sa feuille et les nœuds frères situés sur le chemin
allant de sa feuille à la racine, un par niveau, et recalcule la racine~: $\log_2 n$ hachages pour $n$ comptes, onze
pour deux mille.

\begin{omfigure}
\centering
\begin{tikzpicture}[font=\footnotesize,
  n/.style={draw,thick,rounded corners=2pt,align=center,minimum height=0.75cm,text width=1.35cm},
  e/.style={thick,gray}]
\node[n,draw=omThm,fill=omThm!12] (r) at (5.25,3.3) {racine\\$\Sigma=90$};
\node[n,draw=omExo,fill=omExo!10] (a) at (2.6,2.1) {$h_{12}$\\$\Sigma=35$};
\node[n,draw=omThm,fill=omThm!12] (b) at (7.9,2.1) {$h_{34}$\\$\Sigma=55$};
\node[n] (l1) at (1.3,0.8) {compte 1\\30};
\node[n] (l2) at (3.9,0.8) {compte 2\\5};
\node[n,draw=omExo,fill=omExo!10] (l3) at (6.6,0.8) {compte 3\\40};
\node[n,draw=omThm,fill=omThm!12] (l4) at (9.2,0.8) {compte 4\\15};
\draw[e] (r) -- (a); \draw[e] (r) -- (b); \draw[e] (a) -- (l1); \draw[e] (a) -- (l2);
\draw[e] (b) -- (l3); \draw[e] (b) -- (l4);
\node[anchor=north,align=center,text width=11cm] at (5.25,0.2) {Le compte 4 recalcule sa feuille, puis
$h_{34}$ à partir de la feuille sœur du compte 3, puis la racine à partir du nœud frère $h_{12}$ (ombré
gris)~; il vérifie que chaque somme rencontrée en chemin est positive ou nulle.};
\end{tikzpicture}

\omcaption{Un arbre de Merkle à sommes de quatre comptes. Chaque nœud porte le hachage de ses fils et la somme de
leurs soldes~; la racine engage tous les soldes et leur total. La preuve d'inclusion d'un client est constituée des nœuds frères
de son chemin (en gris pour le compte 4). Schématique.}\label{fig:m3:centralised-exchanges:merkle}
\end{omfigure}

\begin{proposition}[Ce qu'établit une preuve d'inclusion]\label{prop:m3:centralised-exchanges:inclusion}
Soit la fonction de hachage résistante aux collisions. (i) Si la feuille d'un client et son chemin recalculent la
racine publiée, son solde est une feuille de l'arbre engagé et est compté dans son total. (ii) Si de plus chaque
solde de feuille est positif ou nul, le total publié est au moins égal à la somme des soldes de tous les clients dont
les preuves sont valides. Sans (ii), un arbre contenant une feuille négative passe la vérification pour chaque client alors que son
total sous-estime les passifs du montant de cette feuille.
\end{proposition}

\begin{proof}
(i) Une feuille ou un chemin différent aboutissant à la même racine donnerait deux entrées distinctes ayant le même hachage à
un niveau donné. (ii) La somme de la racine est la somme de toutes les feuilles~; retirer les feuilles positives ou nulles des clients
qui n'ont pas vérifié ne peut que la diminuer. Pour la dernière affirmation, une feuille négative ne modifie le hachage d'aucune autre feuille, si bien
que tout chemin honnête recalcule encore la racine, dont la somme est diminuée du montant négatif.
\end{proof}

La feuille négative n'est pas une curiosité~: c'est le moyen le moins coûteux, pour une plateforme insolvable, de réduire ses
passifs publiés pour les faire coïncider avec ses réserves. Elle est détectée lorsque les vérificateurs des clients contrôlent que chaque somme
de leur chemin est positive ou nulle, car la feuille négative rend négatives les sommes de ses ancêtres tant qu'un sous-arbre
détient davantage de soldes réels qu'il n'en dissimule. Le tutoriel construit un arbre de 1\,023 clients et un faux
compte qui dissimule 60~\% des passifs~: les 1\,023 preuves passent un contrôle portant uniquement sur les hachages, et les 1\,023 échouent
à un contrôle des sommes. Les schémas qui ne veulent pas révéler les sommes des nœuds frères recourent à des preuves à divulgation nulle de connaissance sur des intervalles pour montrer
que chaque feuille est positive ou nulle. L'omission de comptes est plus difficile à détecter~: elle ne se révèle que lorsqu'un client omis
vérifie, si bien qu'une plateforme qui omet les comptes inactifs les moins susceptibles de contrôler est repérée avec une faible
probabilité.

\begin{dated}{2026-09}{Une preuve de réserves publiée}\label{dat:m3:centralised-exchanges:kraken}
Kraken décrit sa \omterm{def:m3:centralised-exchanges:por}{preuve de réserves} comme un examen mené par un expert-comptable indépendant~: un instantané anonymisé des
soldes clients agrégés dans un arbre de Merkle, chaque feuille étant les 16 premiers caractères d'un hachage SHA-256 d'un identifiant
d'enregistrement et des soldes, et une comparaison de ces soldes avec les avoirs on-chain dont le contrôle
l'expert-comptable vérifie par des signatures. Les clients peuvent vérifier eux-mêmes leur inclusion dans l'arbre. Kraken énumère
ce que l'examen ne prouve pas~: la possession exclusive des \omterm{def:m3:blockchains-for-traders:key}{clés privées}, l'absence de passifs cachés
ni de fonds empruntés pour réussir l'examen, ni les actifs hors chaîne, ni tout ce qui suit l'instantané~; les actifs
couverts étaient le BTC, l'ETH, le SOL, l'USDC, l'USDT et le XRP.
\end{dated}

Une \omterm{def:m3:centralised-exchanges:por}{preuve de réserves} est donc un instantané d'un côté du bilan, assorti d'un contrôle partiel de l'autre.
Elle ne peut pas montrer les passifs que la plateforme doit à des prêteurs plutôt qu'à ses clients, les actifs prêtés le
lendemain, ni les clés détenues par un tiers. Elle constitue un élément utile contre la fraude la plus grossière, et elle n'aurait pas
révélé le privilège d'emprunt de FTX, qui résidait dans les paramètres du moteur d'appariement, et non
dans les soldes.

\section{Règlement hors plateforme}

Si le risque est que la plateforme détienne le collatéral, le remède est qu'elle ne le détienne pas.

\begin{definition}[Règlement hors plateforme]\label{def:m3:centralised-exchanges:oes}
\emph{règlement hors plateforme}\index{règlement hors plateforme} est un dispositif dans lequel le collatéral d'une société de trading
reste confié à un dépositaire indépendant, la plateforme crédite la société d'un solde de trading en miroir,
et les gains et pertes de son trading sont réglés entre le dépositaire et la plateforme à intervalles fixes ou à la demande, de sorte que
seul le montant non réglé soit exposé au risque pour la plateforme.
\end{definition}

Ce dispositif transfère l'exposition de la société du bilan de la plateforme vers celui du dépositaire, lequel est
choisi pour sa ségrégation, ses contrôles et sa supervision, et limite ce qu'une défaillance de plateforme peut emporter à
une seule période de règlement en gains nets. Il permet aussi à une société de négocier sur plusieurs plateformes avec un même pool de collatéral,
avec une compensation entre elles. Il n'est pas exempt de risques~: le dépositaire devient la concentration, la plateforme accorde
un crédit à la société entre deux règlements et l'intègre dans ses limites, et la société a besoin que ses
positions, ses soldes en miroir et les règlements se réconcilient à chaque cycle.

\begin{dated}{2026-09}{Un réseau de règlement hors plateforme}\label{dat:m3:centralised-exchanges:clearloop}
Copper, dépositaire, décrit son réseau ClearLoop comme « \omterm{def:m3:centralised-exchanges:oes}{règlement hors plateforme} » assorti d'un « règlement net
unique » entre les plateformes connectées, qu'elle énumère comme incluant OKX, Coinbase, Bybit, Gate.io, Deribit,
Bitget, Bitfinex et Kraken~; les actifs sont conservés sous la garde de Copper. D'autres dépositaires et certaines plateformes proposent
des dispositifs similaires~; les conditions, la fréquence des règlements et les plateformes connectées relèvent des prestataires et changent.
\end{dated}

\begin{omfigure}
\centering
\begin{tikzpicture}[font=\small,
  p/.style={draw,thick,rounded corners=2pt,align=center,minimum height=1cm,text width=2.7cm},
  arr/.style={-{Stealth[length=2.2mm]},thick}]
\node[p,draw=omDef,fill=omDef!8] (f) at (0,0) {société de trading};
\node[p,draw=omProp,fill=omProp!8] (c) at (5,0) {dépositaire\\indépendant\\(collatéral)};
\node[p,draw=omThm,fill=omThm!8] (x1) at (10,0.9) {plateforme A\\(solde en miroir)};
\node[p,draw=omThm,fill=omThm!8] (x2) at (10,-0.9) {plateforme B\\(solde en miroir)};
\draw[arr] (f) -- node[above,font=\footnotesize] {dépôts} (c);
\draw[arr] (c) -- node[above,sloped,font=\footnotesize] {ligne de crédit} (x1);
\draw[arr] (c) -- node[below,sloped,font=\footnotesize] {ligne de crédit} (x2);
\draw[arr,dashed] (f.north) to[bend left=25] node[above,font=\footnotesize] {ordres} (x1.north);
\draw[arr,dashed] (f.south) to[bend right=25] node[below,font=\footnotesize] {ordres} (x2.south);
\end{tikzpicture}

\omcaption{\omterm{def:m3:centralised-exchanges:oes}{Règlement hors plateforme}~: le collatéral de la société reste confié à un dépositaire, les plateformes créditent
un solde de négociation en miroir, et les gains et pertes nets sont réglés entre le dépositaire et les plateformes à heures fixes.
La défaillance d'une plateforme ne peut emporter que le montant non réglé. Schéma.}\label{fig:m3:centralised-exchanges:oes}
\end{omfigure}

\section{Tutoriel~: vérifier une preuve de réserves}\label{tut:m3:centralised-exchanges:por}

\begin{tutorial}
\textbf{Objectif.} Construire un arbre de sommes de Merkle des soldes des clients, produire et vérifier des preuves d'inclusion, et
construire le contre-exemple de la \cref{prop:m3:centralised-exchanges:inclusion}. \textbf{État final~:} les
chiffres de la section consacrée à la \omterm{def:m3:centralised-exchanges:por}{preuve de réserves}.
\begin{tutsteps}
\item \textbf{L'arbre.} Les feuilles engagent le compte, le nonce et le solde~; les nœuds parents engagent leurs enfants et leur somme.
\omcode{code/markets-3/15-centralised-exchanges/python/m3_cex.py}{14}{52}{Un arbre de sommes de Merkle, des preuves d'inclusion et leur vérification.}\label{lst:m3:centralised-exchanges:tree}
\item \textbf{Le contre-exemple.} Insérer un compte dont le solde est négatif et égal à 60~\% des
passifs, puis compter les clients dont les preuves passent avec et sans la vérification des sommes~:
\texttt{fake\_negative(customers())}.
\item \textbf{Le déficit.} \texttt{ftx\_coverage()} lit le tableau des débiteurs et renvoie les parts du
\cref{tab:m3:centralised-exchanges:shortfall}~; \texttt{fig\_cex.py} écrit les données du graphique du chapitre.
\end{tutsteps}
\textbf{Pour aller plus loin.} Masquer seulement 2~\% des passifs et compter les clients qui le détectent encore~;
placer la feuille fictive à côté du plus gros compte~; remplacer les sommes des nœuds frères par des engagements, et se demander
ce qu'un client peut encore vérifier.
\end{tutorial}

\section{Construction~: le régulateur de limite de débit}\label{bld:m3:centralised-exchanges:ratelimit}

\begin{build}
\bfield{Rôle} Chaque requête que la petite société envoie à une plateforme passe par un régulateur qui connaît les
limites de la plateforme et refuse, avant l'envoi, toute requête que la plateforme rejetterait~; après un code 429 ou 418, il
s'arrête pendant le temps indiqué par la plateforme. Un système de trading qui apprend les limites à partir des rejets de la plateforme se fait bannir.
\bfield{Interface} \texttt{Rule(kind, interval\_ms, limit)} pour les règles de poids ou de nombre d'ordres~;
\texttt{Governor(rules)}~; \texttt{try\_send(now, weight, is\_order)} renvoie s'il faut envoyer et, sinon,
combien de temps attendre~; \texttt{on\_status(now, status, retry\_after\_ms)}. L'en-tête C++20 \texttt{firm\_ratelimit.hpp}
et la crate Rust \texttt{firm\_ratelimit} implémentent la même interface~; le module Python fait référence.
\bfield{Règles} Fenêtres fixes alignées sur l'époque, comme la plateforme les compte~; une requête n'est envoyée que si chaque règle
dispose de place, et consomme alors toutes les règles~; un régulateur bloqué refuse tout jusqu'à la fin du blocage~; entiers
en millisecondes, sans virgule flottante.
\bfield{Tests d'acceptation} \texttt{code/firm/ratelimit/tests/}, \texttt{cpp/firm\_ratelimit\_test.cpp} et les
tests de la crate Rust~: le temps d'attente d'une requête refusée jusqu'à la fenêtre suivante~; les règles d'ordres et de poids comptées
séparément~; la temporisation après les codes 429 et 418.
\bfield{Pour aller plus loin} Fenêtres glissantes~; poids par point de terminaison lus dans la réponse d'information sur l'échange de la plateforme~;
réconciliation avec les en-têtes d'usage de la plateforme~; un émetteur de cotations à regroupement au-dessus du régulateur.
\end{build}

\omcode{code/firm/ratelimit/cpp/firm_ratelimit.hpp}{41}{60}{La décision du régulateur C++20 et sa réaction aux codes d'état de la plateforme.}\label{lst:m3:centralised-exchanges:governor}

\begin{omsources}
\item US Securities and Exchange Commission, ``SEC Charges Samuel Bankman-Fried with Defrauding Investors
in Crypto Asset Trading Platform FTX'', communiqué de presse 2022-219, 13 décembre 2022.
\item FTX Debtors, ``Preliminary Analysis of Shortfalls'', présentation au Official Committee of Unsecured
Creditors, In re FTX Trading Ltd., No.\ 22-11068 (JTD), Bankr.\ D.\ Del., Doc 792-1, 2 mars 2023.
\item J.~J.~Ray III, \emph{First Interim Report to the Independent Directors on Control Failures
at the FTX Exchanges}, Doc 1242-1, 9 avril 2023~; FTX Debtors, plainte contre les fondateurs, Doc 1886, 20 juillet 2023.
\item Securities and Exchange Commission v.\ Bankman-Fried, plainte, S.D.N.Y., 13 décembre 2022~; FTX Debtors, communiqué de
presse sur la confirmation du plan, 7 octobre 2024~; NPR, reportage sur le dépôt de bilan de Mt.~Gox, 28 février 2014.
\item Commodity Futures Trading Commission, communiqué de presse 8938-24, 8 août 2024~; UPI, compte rendu de la
condamnation, 28 mars 2024.
\item Binance, Spot API documentation (limits, enums) et fee schedule~; Kraken, fee schedule et pages de preuve
de réserves~; Copper, ClearLoop product page~; site du trustee de la procédure de réhabilitation de MtGox. Tous consultés en
septembre 2026.
\end{omsources}

\section{Exercices}

\begin{exercise}[$\star$]\label{exo:m3:centralised-exchanges:1}
Nommer les quatre rôles que jouent des institutions distinctes dans une transaction à terme régulée et que
qu'une \omterm{def:m3:centralised-exchanges:cex}{plateforme centralisée} assure seule.
\end{exercise}

\begin{exercise}[$\star$]\label{exo:m3:centralised-exchanges:2}
Une limite est de 1\,200 unités de poids par minute en fenêtres fixes. Quel est le maximum qu'un client peut envoyer en
deux secondes quelconques~? En une minute quelconque~?
\end{exercise}

\begin{exercise}[$\star$]\label{exo:m3:centralised-exchanges:3}
Pourquoi une \omterm{def:m3:centralised-exchanges:subaccount}{clé d'API} utilisée par une stratégie de trading ne devrait-elle pas porter de droits de retrait~?
\end{exercise}

\begin{exercise}[$\star\star$]\label{exo:m3:centralised-exchanges:4}
Une stratégie échange 50~millions de dollars par jour, moitié en maker et moitié en taker, au palier régulier de Binance puis
au palier VIP~9 (\cref{dat:m3:centralised-exchanges:fees}). Combien paie-t-elle en frais par jour dans chaque cas~?
\end{exercise}

\begin{exercise}[$\star\star$]\label{exo:m3:centralised-exchanges:5}
Une plateforme compte 2\,000\,000 comptes. Combien de hachages un client calcule-t-il pour vérifier son inclusion~? Si
la plateforme omet 100 comptes et que chaque client vérifie avec une probabilité de 1~\%, quelle est la probabilité
qu'au moins une omission soit détectée~?
\end{exercise}

\begin{exercise}[$\star\star$]\label{exo:m3:centralised-exchanges:6}
D'après le \cref{tab:m3:centralised-exchanges:shortfall}, quelle part des créances en espèces et en \omterm{def:m3:blockchains-for-traders:stable}{stablecoins} la plateforme aurait-elle pu payer
avec les espèces, les \omterm{def:m3:blockchains-for-traders:stable}{stablecoins} et les créances de cette ligne (270 et 310 millions de dollars)~?
\end{exercise}

\begin{exercise}[$\star\star\star$]\label{exo:m3:centralised-exchanges:7}
\emph{Codage.} Avec \texttt{fake\_negative}, masquer 2~\% des passifs au lieu de 60~\%. Combien parmi les
1\,023 clients le détectent lorsqu'ils vérifient les sommes~? Expliquer pourquoi ils sont moins nombreux.
\end{exercise}

\begin{exercise}[$\star\star\star$]\label{exo:m3:centralised-exchanges:8}
\emph{Trouver la faille.} « La plateforme a publié le mois dernier une \omterm{def:m3:centralised-exchanges:por}{preuve de réserves} montrant des actifs supérieurs aux soldes des clients,
donc nos 40~millions de dollars déposés là-bas sont en sûreté. »
\end{exercise}

\section{Problème~: le trou du bilan}

\begin{problem}[{Problème du week-end --- ce qu'un client de FTX.com devait recevoir et ce qui était disponible}]\label{pb:m3:centralised-exchanges:1}
Utiliser le \cref{tab:m3:centralised-exchanges:shortfall}, \cref{fig:m3:centralised-exchanges:flows} et les
chiffres des débiteurs cités dans le texte. Tous les montants sont en millions de dollars, aux prix à la date du dépôt de bilan.

\medskip\noindent\textbf{Partie I --- La ruée.}
\begin{enumerate}
\item Quel montant les clients ont-ils retiré, net, de FTX.com du 1 au 11 novembre 2022~?
\item Quel jour le retrait net a-t-il été le plus important, et de quel montant~?
\item Que firent les parties liées pendant ces mêmes jours, et pourquoi cela pourrait-il importer aux clients~?
\item Pourquoi les retraits se sont-ils presque arrêtés à partir du 9 novembre~?
\item Qui a été intégralement remboursé lors de cette ruée, et pourquoi~?
\end{enumerate}

\medskip\noindent\textbf{Partie II --- Le déficit.}
\begin{enumerate}[resume]
\item Quelle part des créances en \omterm{def:m3:blockchains-for-traders:key}{jetons} liquides était couverte par les \omterm{def:m3:blockchains-for-traders:key}{jetons} liquides détenus~?
\item Quelle part l'était avec les créances sur ces \omterm{def:m3:blockchains-for-traders:key}{jetons} (383~millions de dollars)~?
\item Quelle part des créances en bitcoins était couverte~?
\item Quelle part de l'ensemble des créances était couverte par l'ensemble des actifs et des créances, et quel était le déficit~?
\item Pourquoi cette dernière part est-elle trompeuse~?
\end{enumerate}

\medskip\noindent\textbf{Partie III --- Ce qui aurait permis de le voir.}
\begin{enumerate}[resume]
\item Une \omterm{def:m3:centralised-exchanges:por}{preuve de passif} par arbre de Merkle aurait-elle révélé le déficit~?
\item Une \omterm{def:m3:centralised-exchanges:por}{preuve de réserves}, comparant les avoirs on-chain au total, l'aurait-elle révélé~?
\item L'un ou l'autre aurait-il révélé le privilège d'emprunt d'Alameda~?
\item Comment le \omterm{def:m3:centralised-exchanges:oes}{règlement hors plateforme} aurait-il modifié la perte d'une société~?
\item Quelle aurait dû être la limite d'une société sur une plateforme, et mesurée comment~?
\end{enumerate}

\medskip\noindent\textbf{Partie IV --- Jugement.}
\begin{enumerate}[resume]
\item Pourquoi les créances ont-elles été fixées en dollars aux cours de la date du dépôt de bilan, et qui y a gagné ou perdu~?
\item En quoi le solde d'un client sur une \omterm{def:m3:centralised-exchanges:cex}{plateforme centralisée} diffère-t-il d'un compte chez un courtier
avec des actifs clients ségrégués~?
\item Quels signes, avant le 6 novembre, un responsable des risques aurait-il pu exploiter~?
\item Énoncer le \emph{résultat nommé}~: la part des créances des clients de FTX.com en \omterm{def:m3:blockchains-for-traders:key}{jetons} liquides couverte par les
\omterm{def:m3:blockchains-for-traders:key}{jetons} liquides détenus à la date du dépôt de bilan, sans et avec les créances.
\item En une phrase~: qu'est-ce qu'un solde sur une \omterm{def:m3:centralised-exchanges:cex}{plateforme centralisée}~?
\end{enumerate}
\end{problem}

\section{Questions d'entretien}

\begin{interviewq}[$\star$ \iqroles{trader}]\label{iq:m3:centralised-exchanges:1}
Qu'est-ce qu'un \omterm{def:m3:centralised-exchanges:postonly}{ordre post-only}, et quand l'utiliseriez-vous~?
\end{interviewq}

\begin{interviewq}[$\star$ \iqroles{risque}]\label{iq:m3:centralised-exchanges:2}
Comment limiteriez-vous l'exposition d'une société de trading à la défaillance d'une plateforme de crypto-actifs~?
\end{interviewq}

\begin{interviewq}[$\star\star$ \iqroles{développeur}]\label{iq:m3:centralised-exchanges:3}
Votre système reçoit une erreur HTTP 429 d'une plateforme plusieurs fois par heure. Que faites-vous~?
\end{interviewq}

\begin{interviewq}[$\star\star$ \iqroles{développeur, chercheur}]\label{iq:m3:centralised-exchanges:4}
Expliquez comment un client vérifie son inclusion dans la \omterm{def:m3:centralised-exchanges:por}{preuve de passif} d'une plateforme, et
comment elle pourrait la truquer.
\end{interviewq}

\begin{interviewq}[$\star\star$ \iqroles{trader, risque}]\label{iq:m3:centralised-exchanges:5}
Que ne prouve pas une \omterm{def:m3:centralised-exchanges:por}{preuve de réserves}~?
\end{interviewq}

\begin{interviewq}[$\star\star\star$ \iqroles{développeur}]\label{iq:m3:centralised-exchanges:6}
Concevez le composant qui envoie les mises à jour de cotations pour 200 instruments à une plateforme
qui autorise 100 ordres par 10 secondes et par compte.
\end{interviewq}
